\section{Introduction to Discrete Mathematics}\label{dm:section:1.idm}

\subsection{Arrangements and Permutations}

\subsubsection{Arrangements}

\definition{13.1.1.1}{An \underline{arrangement} is a sequence of items in which the order of the items in the sequence matters.}

\separate{7}\definition{13.1.1.2}{The \underline{length} of an arrangement is the number of items in the sequence.}

\separate{7}\standardproblem{1}{How many arrangements of length \(k\) can be made from elements of a set \(S\) of size \(|S|=n\)?}
{\(n^k\) as arrangement has \(k\) places, each of which can be occupied by any one of \(n\) entries.}

\separate{7}\example{13.1.1.1}{How many four-letter \say{words} can be made from the standard A-Z alphabet?

\separate{3}\textit{Solution}: \(26^4 = 456,796\) words.}

\separate{7}\standardproblem{2}{How many arrangements of length \(k\) with no repeats can be made from elements of a set \(S\) of size \(|S|=n\)?}
{\begin{equation} 
    n(n-1)(n-2)...(n-k+1) = \frac{n!}{(n-k)!},\qquad k \leq n    
\end{equation}
As we look at the available choices for each successive item in the list, the number of options is reduced by one each time.}

\separate{7}\definition{13.1.1.3}{The \underline{multiplication principle} states that, for an arrangement of length \(k\), if for \(i \in \{1,\;2,\;...,\;k\}\) there are \(n_i\) choices for the \(n^{th}\) item, regardless of previous choices, then the number of possible arrangements is \(n_1n_2...n_k\).}

\subsubsection{Permutations}

\definition{13.1.2.1}{A \underline{permutation} on a finite set \(S\) is an arrangement of length \(|S|\) that uses each item from \(S\) exactly once.}

\separate{7}\proposition{13.1.2.1}{If \(|S|=n\), then there are \(n!\) permutations on \(S\).}
{This is \hyperref[prob:2]{Standard Problem 2} with \(n=k\).}

\separate{7}\definition{13.1.2.2}{We generalise this concept as, for a non-negative integer \(r\), an \underline{\(r\)-permutation} on \(S\) is an arrangement of \(r\) distinct elements from \(S\). We write \(P(n,r)\) for the number of \(r\)-permutations on a set of size \(n\), which, by \hyperref[prob:2]{Standard Problem 2}, we have,
\begin{equation}
    P(n,r) = n(n-1)...(n-r+1) = \frac{n!}{(n-r)!},\qquad 0 \leq r \leq n
\end{equation}}

\subsubsection{Counting Examples}

\definition{13.1.3.1}{The \underline{addition principle} states that, for two disjoint, finite sets \(A\) and \(B\). Then,
\begin{equation}
    A \cap B = \emptyset \qquad |A \cup B| = |A| + |B|
\end{equation}
More generally, if \(A_1,\;A_2,\;...,\;A_n\) are pairwise disjoint, then
\begin{equation}
    |A_1 \cup A_2 \cup ... \cup A_n| = |A_1| + |A_2| + ... + |A_n|
\end{equation}}

\subsection{Combinations}

\subsection{Arrangements and Combinations with Repetitions}

\section{Combinatorics}\label{dm:section:2.com}

\subsection{The Three Principles}

\subsection{Recurrence Relations}

\subsection{Generating Functions}

\section{Graph Theory}\label{dm:section:3.gt}

\subsection{What is a Graph?}

\subsection{Eulerian Graphs}

\subsection{Planar Graphs}